\documentclass[11pt,a4paper]{article}
\usepackage[margin=1in]{geometry}
\usepackage{booktabs}
\usepackage{array}
\usepackage{tabularx}
\usepackage{hyperref}
\usepackage{enumitem}
\usepackage{titlesec}
\usepackage{xcolor}
\usepackage{fancyhdr}

\hypersetup{
  colorlinks=true,
  linkcolor=black,
  urlcolor=teal!60!black
}

\pagestyle{fancy}
\fancyhf{}
\fancyhead[L]{\small Project Summary}
\fancyhead[R]{\small Supaporn Klabklaydee}
\fancyfoot[C]{\thepage}
\renewcommand{\headrulewidth}{0.4pt}

\titleformat{\section}{\large\bfseries}{\thesection}{0.75em}{}
\titleformat{\subsection}{\normalsize\bfseries}{\thesubsection}{0.75em}{}

\begin{document}

\begin{center}
  {\LARGE\bfseries Project Summary}\\[0.6em]
  {\large Sample-Efficiency of Quantum Fidelity Kernels\\
  for Molecular Property Classification:\\[0.25em]
  A Learning-Curve Benchmark Against Classical Baselines}\\[1em]
  {\normalsize Supaporn Klabklaydee}\\
  {\small Institute of Science Tokyo}\\[0.4em]
  {\footnotesize A negative / null-result study with multi-seed, scaffold-split learning curves.}
\end{center}

\vspace{0.6em}
\noindent\rule{\textwidth}{0.6pt}
\vspace{0.8em}

\section*{Core Question}
Does a quantum fidelity kernel (simulated on a classical computer) show a real \emph{sample-efficiency} advantage over classical machine-learning methods when labeled molecular data is scarce---the low-data regime that matters most in drug discovery and toxicology, where labels are expensive?

\section*{Method}
The study benchmarks a PennyLane simulator-based quantum fidelity kernel SVM
(8-qubit ``Circuit~B,'' using angle encoding + ring CNOT entanglement on PCA-8 molecular features)
against three classical baselines---Random Forest, RBF-SVM, and AttentiveFP graph neural network---
across four MoleculeNet binary classification tasks:
\begin{itemize}[leftmargin=1.4em,itemsep=0.2em]
  \item \textbf{BBBP} --- blood--brain barrier permeability
  \item \textbf{BACE} --- enzyme inhibition
  \item \textbf{HIV} --- replication inhibition
  \item \textbf{Tox21 NR-AR} --- toxicity
\end{itemize}

\paragraph{Key design rigor.}
Training sizes were swept from $N=50$ up to the full dataset (7 sizes),
with Bemis--Murcko scaffold splitting (to prevent structural leakage between train and test)
and 10 random seeds---\textbf{280 total model--condition pairs}.
A pre-registered hypothesis was tested: that the quantum kernel would show a statistically significant advantage over the GNN specifically at $N<500$ (paired Wilcoxon, $p<0.05$).

\section*{Key Results}

\begin{center}
\renewcommand{\arraystretch}{1.25}
\begin{tabularx}{\textwidth}{@{}>{\bfseries}l X X@{}}
\toprule
Dataset & Quantum vs.\ GNN advantage & Quantum vs.\ best classical model \\
\midrule
BBBP & Significant at $N=300$, $N=500$ & Never exceeded RF / RBF-SVM \\
BACE & No significant advantage & Quantum near chance-level ($\sim$0.50 AUC) throughout \\
HIV & Significant at $N=50$--$200$ & Never exceeded RF / RBF-SVM \\
Tox21 & Mean crossover at $N=50$, not significant ($p=0.23$) & Not robust; RBF-SVM later dominated \\
\bottomrule
\end{tabularx}
\end{center}

\vspace{0.5em}
The quantum kernel's AUC stayed flat in a narrow \textbf{0.46--0.53} band on three of four tasks---essentially chance-level---while Random Forest and RBF-SVM generally dominated overall, with gaps of \textbf{0.12--0.34 AUC} points at larger $N$.

\section*{Why the ``Significant'' Wins Are Not What They Look Like}
The paper is unusually candid: the quantum-over-GNN significant results largely reflect \textbf{GNN underfitting at small $N$}, not quantum kernel strength.
On HIV at $N=50$, the GNN scored only \textbf{0.340} while Random Forest already hit \textbf{0.552}---the quantum kernel's \textbf{0.491} beat the GNN but still trailed the simpler classical baselines.

\section*{Explanation Offered}
The authors attribute the quantum kernel's plateau to \textbf{expressivity limits}:
with 8~qubits, the fidelity kernel's rank is bounded by a 256-dimensional Hilbert space and constrained by the fixed ring-CNOT circuit structure,
unlike the RBF kernel or Morgan fingerprints, which operate in much higher-dimensional, more flexible spaces.
This points to a \emph{capacity-saturated} regime rather than overfitting.

\section*{Practical Trade-offs}
Computing the quantum kernel's Gram matrix required $O(N^2)$ circuit evaluations,
taking \textbf{hours per dataset} on a simulator, versus \textbf{seconds to minutes} for classical Random Forest and RBF-SVM on the same hardware---
a substantial computational cost for no net performance benefit.

\section*{Conclusion}
This is presented as the first systematic scaffold-split learning-curve comparison of a quantum fidelity kernel against modern classical baselines with proper multi-seed statistics.

\textbf{Bottom line:} the quantum kernel is not globally competitive with the best classical models on any of the four benchmarks,
though it does show narrow, likely GNN-underfitting-driven advantages in a few low-data corners (BBBP, HIV).
The authors are transparent about limitations---no Bonferroni correction across the many comparisons, simulator-only evaluation, and fixed (non-trainable) quantum circuit design---
and frame the exploratory ``wins'' as evidence of GNN fragility at low~$N$ rather than proof of quantum advantage.

\vspace{0.6em}
\noindent\fbox{\parbox{0.97\textwidth}{\small
\textbf{Takeaway.}
This is a negative / null-result paper with strong methodological rigor---valuable as a reproducible baseline for the field precisely because it resists overselling a fashionable ``quantum advantage'' narrative.
}}

\end{document}
